%%% FORMATO E IDIOMA %%%
    \documentclass[letterpaper,DIV=15,12pt]{scrartcl}
    \usepackage[spanish,mexico,shorthands=off,es-lcroman]{babel}
    \usepackage{scrlayer-scrpage}
        \clearpairofpagestyles
        \ihead{\footnotesize \textit{Teoría de Conjuntos III}}
        \ohead{\footnotesize \textit{2026-2}}
        \cfoot{\normalfont\thepage}
        \addtokomafont{title}{\bfseries \rmfamily}
        \setlength{\headsep}{5pt}
        \newpairofpagestyles{beginstyle}{
            \clearpairofpagestyles
            \KOMAoptions{headsepline=false}
            \cfoot{\footnotesize \pagemark}
            \cfoot{\normalfont\thepage}
        }
        \addtokomafont{section}{\raggedleft\rmfamily}
        \addtokomafont{subsection}{\raggedleft\rmfamily}
        \addtokomafont{subsubsection}{\raggedleft\rmfamily}
        \setlength{\headsep}{8pt}
        \setlength{\footskip}{20pt}
        \setlength{\textheight}{600pt}
%%% PAQUETES %%%
    \usepackage{ifthen}
    \usepackage[shortlabels]{enumitem}
        \setenumerate[1]{label=\MakeLowercase{\roman*}), ref=\roman*}
        \setenumerate[2]{label=\MakeLowercase{\alph*}), ref=\alph*}
    \usepackage{stix2}
    \usepackage{amsmath}
    \usepackage{amsthm}
	\renewcommand{\qedsymbol}{$\blacksquare$}
	\addto\captionsspanish{\renewcommand{\proofname}{\normalfont\bfseries Demostración}}	
	\newenvironment{solucion}{\renewcommand{\qedsymbol}{$\boxtimes$}\begin{proof}[\normalfont\bfseries Solución]}{\end{proof}}
%%% E J E R C I C I O S %%%
    \newcounter{Ejer}
    \newcounter{Pts}
    \setcounter{Ejer}{1}
    \setcounter{Pts}{1}
    \newcommand{\pts}{}
    \newenvironment{ejercicio}[1]{\noindent
        \ifthenelse{\equal{#1}{1}}{\renewcommand{\pts}{\textbf{(#1 pt)}}}{\renewcommand{\pts}{\textbf{(#1 pts)}}}\textbf{Ej. \theEjer} \pts\stepcounter{Ejer}}{}
%%% C O M A N D O S %%%
    \renewcommand{\emptyset}{\varnothing}
    \newcommand{\tq}{\mid}
%% D O C U M E N T O %%
\begin{document}
\thispagestyle{plain}
    
       \begin{center}
            {\fontsize{30}{60}\rmfamily \textbf{Ejercicio semanal 2}} \\ \vspace{.2cm} Teoría de Conjuntos III, 2026-2
       \end{center}
        \begin{flushright}
            \footnotesize \hfill Profesor: Luis Jesús Turcio Cuevas.\\
            \hfill Ayudante: Hugo Víctor García Martínez.
        \end{flushright}
   
       \textit{\textbf{\textsc{Instrucciones.}} Esta tarea es \textbf{individual} y deberá ser entregada \textbf{presencial y personalmente} el día \textbf{miércoles 10 de febrero} al inicio la clase.}\vspace{.3cm}

    \begin{ejercicio}{10}
        Sea $\mathbb{P}=(P,\leq)$ un orden parcial y $G \subseteq P$ un filtro de $\mathbb{P}$. Pruebe que $G$ es un filtro genérico si y sólo si para toda anticadena maximal $A$ de $\mathbb{P}$, $A \cap G \neq \emptyset$.
    \end{ejercicio}
    \begin{proof}
        Supóngase primero que $G$ es filtro genérico y sea $A$ una anticadena maximal de $\mathbb{P}$. Entonces $D=\downarrow A = \{ p \in P \tq \exists a \in A \, (p \leq a) \}$ y existen $g \in G$ y $a \in A$ tales que $g \leq a$. Puesto que $G$ es filtro, es cerrado bajo cotas superiores, por lo que $g \in A$ y $G \cap A\neq \emptyset$.

        Ahora, supóngase que $G$ interseca a cada anticadena maximal de $\mathbb{P}$ y sea $D \subseteq P$ denso. Como $D$ es denso, contiene una anticadena maximal (demostrado en clase), por lo que $G$ interseca a tal anticadena, y con ello $G \cap D \neq \emptyset$.
    \end{proof}
\end{document}